\documentclass[11pt]{article}
\usepackage[T1]{fontenc}
\usepackage{lmodern}
\usepackage[margin=1in]{geometry}
\usepackage{amsmath,amssymb,amsthm,booktabs}
\usepackage{xcolor}
\usepackage[colorlinks,allcolors=blue!50!black]{hyperref}

\newtheorem{theorem}{Theorem}
\newtheorem{lemma}{Lemma}
\newcommand{\disc}{\operatorname{Disc}}
\newcommand{\real}{\operatorname{Re}}
\newcommand{\certbase}{https://github.com/harrykane8816627373-cell/mordell-research}
\newcommand{\certrev}{cf631a15c2d62b61b6f40b5e16092a4c5a189978}
\newcommand{\certfile}[2]{\href{\certbase/blob/\certrev/#1}{#2}}

\title{The Sharp Ten-Point Bound for Mordell's Distance Product}
\author{Fengcan Jiang}
\date{26 November 2026}
\hypersetup{pdftitle={The Sharp Ten-Point Bound for Mordell's Distance Product},
  pdfauthor={Fengcan Jiang},
  pdfsubject={Sharp distance-product inequality and equality classification}}

\begin{document}
\maketitle

\begin{abstract}
We determine the largest product of the squared pairwise distances among ten
points in the plane with sum of squared radii equal to ten. The maximum is
an explicitly specified algebraic number, approximately
\(5.3823025895657\times 10^{10}\). Up to rotation and relabeling, the unique
maximizer consists of an inner antipodal pair and two concentric rectangles.
The proof uses a two-point majorant indexed by cyclic separation, with
exact interpolation at fifteen pair types of the candidate. An
integer-arithmetic verifier covers the six necessary pair domains and
certifies the bound with zero slack.
\end{abstract}

\section{The sharp value}

For \(z=(z_1,\ldots,z_n)\in\mathbb C^n\), define
\[
 Q(z)=\sum_{i=1}^n |z_i|^2,\qquad
 \Delta(z)=\prod_{i<j}|z_i-z_j|^2,\qquad
 M_n=\max_{Q(z)=n}\Delta(z).
\]
Ji, Xu, and Yu \cite{jxy} give the sharp values through nine points and
an explicit ten-point configuration exceeding the product of a regular
nonagon together with its center. We give the sharp value and equality
classification at ten points. The proof uses their eight- and nine-point
bounds in its initial restrictions on a global maximizer. The proof at
twelve points \cite{n12} also uses an exact contact kernel, and here the
cyclic order of the eight outer points supplies further zero-sum terms.

Define \(a,b,c,x,y\) as the unique real root of the system
\begin{equation}
 E_k=2\sum_{j=0}^kS_j S_{k-j}-(2k+1)S_k-10T_k=0,
 \qquad 0\le k\le4,                                           \label{eq:root}
\end{equation}
in the closed box with center
\(2^{-20}(118037,1084060,1216217,-797427,810046)\) and coordinate
half-width \(2^{-20}\). Here \(A_k=a^k\), and
\[
\begin{aligned}
 B_0&=2, & B_1&=2x, &
 B_k&=2xB_{k-1}-b^2 B_{k-2},\\
 C_0&=2, & C_1&=2y, &
 C_k&=2yC_{k-1}-c^2 C_{k-2}\qquad (k\ge2),
\end{aligned}
\]
with \(S_k=A_k+B_k+C_k\) and \(T_k=aA_k+bB_k+cC_k\). Put
\[
 P(w)=(w-a)(w^2-2xw+b^2)(w^2-2yw+c^2).
\]
Let \(z^*\) be the unordered configuration
\begin{equation}
 \{\pm\sqrt a\}
 \ \cup\
 \left\{\pm\sqrt{\frac{b+x}{2}}\ \pm i\sqrt{\frac{b-x}{2}}\right\}
 \ \cup\
 \left\{\pm\sqrt{\frac{c+y}{2}}\ \pm i\sqrt{\frac{c-y}{2}}\right\},
                                                               \label{eq:config}
\end{equation}
where the two signs in each rectangle are independent.

\begin{theorem}\label{thm:main}
Under \(Q(z)=10\), the sharp bound is
\[
 M_{10}=\left(\frac{10}{9}\right)^{45}
         4^5 ab^2c^2\disc(P)^2.
\]
It lies in the interval
\([53823025895.656995698750,\,53823025895.656995698751]\).
Every maximizing configuration is a rotation and relabeling of
\(\sqrt{10/9}\,z^*\).
\end{theorem}

\section{Normalization and the candidate}

We work at \(Q=9\). The product has degree \(90\) under dilations of the
coordinates, so the conversion to \(Q=10\) multiplies it by
\((10/9)^{45}\). A global maximum exists on the compact mass sphere and
has distinct entries, since some feasible configurations have positive
product. Each such maximizer also globally maximizes the unconstrained
function
\[
 L(z)=\log\Delta(z)-5Q(z).
\]
Indeed, along the dilates \(sz\) of any point on the mass-nine sphere,
\(L(sz)=\log\Delta(z)+90\log s-45s^2\), with its unique maximum at \(s=1\).
The stationarity equations are
\begin{equation}
 \sum_{j\ne i}\frac{1}{z_i-z_j}=5\bar z_i,\qquad 1\le i\le10.
                                                               \label{eq:force}
\end{equation}
The real Hessian of \(L\) is negative semidefinite at each global
maximizer.

The five roots \(w_i\) of \(P\) are the squares of representatives of the
five antipodal pairs in \eqref{eq:config}, with \(r_i=|w_i|\).
Multiplication of the residual in \eqref{eq:force} by a representative
of the \(i\)-th pair gives
\[
 {\cal E}_i=\tfrac12+2w_i\sum_{j\ne i}(w_i-w_j)^{-1}-5r_i.
\]
The two terms in the sum of forces from another antipodal pair combine
to \(2z_i/(w_i-w_j)\). Expanding the divided difference of \(w^{k+1}\)
gives \(E_k=2\sum_i w_i^k{\cal E}_i\). The \(w_i\) are distinct in the
root box, so its Vandermonde matrix is invertible and \eqref{eq:root}
implies the full force equations. The equation \(E_0=0\) gives
\(a+2b+2c=9/2\). Thus \(z^*\) is feasible and stationary. The exact
root check \cite{data} applies the contraction theorem to
\(v\mapsto v-C E(v)\), where the preconditioner \(C\) is an invertible
rational matrix. The image lies strictly inside the box and the
verified infinity norm of \(I-C\,DE\) is less than \(21/10000\).

An antipodal pair with squared representative \(w_i\) contributes
\(4|w_i|\) to \(\Delta(z^*)\). The four distances squared across two
antipodal pairs have product \(|w_i-w_j|^4\). Multiplying gives
\(\Delta(z^*)=4^5 ab^2c^2\disc(P)^2\). Write \(B=\log\Delta(z^*)\).

\section{Necessary pair domains and cyclic order}\label{sec:domains}

For any pair \(z,w\), use
\[
 u=|z|^2,\qquad v=|w|^2,\qquad t=\real(z\bar w),\qquad
 d=u+v-2t.
\]
The Gram condition is \(t^2\le uv\). The following necessary conditions
hold at all global maximizers at \(Q=9\). Each entry links the
analytic reduction and exact verifier inputs at the fixed revision
\cite{data}.

\begin{center}
\begin{tabular}{ll}
\toprule
Necessary condition & Proof and certificate\\
\midrule
\(d>2/5\) for every pair
 & \certfile{maximizer-domain/README.md}{Initial domain}\\
Two radii squared in \(I=(21/250,159/1000)\)
 & \certfile{ten-point-angular-radii/README.md}{Radial population and bounds}\\
Eight radii squared in \(O=(89/100,1271/1000)\) &\\
\(43/100<d<53/100,\ t/\sqrt{uv}<-87/100\) on the inner edge
 & \certfile{ten-point-inner-edge/README.md}{Inner edge}\\
\(59/100<\cos\gamma<39/50\) on adjacent outer edges
 & \certfile{ten-point-outer-gap/README.md}{Adjacent outer gaps}\\
\(t/\sqrt{uv}<-31/32\) at cyclic separation four
 & \certfile{ten-point-tight-pairing/README.md}{Near-antipodal pairing}\\
\bottomrule
\end{tabular}
\end{center}

The initial reduction uses deletion and recentering, applying the
sharp bounds for eight and nine points to the remaining points. The
radial and angular checks use a coarser certified upper kernel to bound
a sum of pair deficits. They combine the deficit budget with exact
force identities along each row of pairs, population counts and the
fact that the eight successive outer gaps sum to \(2\pi\). The cited
proofs give the implications from these bounds to the strict
inequalities in the table.

Set \(\alpha=\arccos(39/50)\) and \(\beta=\arccos(59/100)\).
Successive outer directions are separated by angles in
\((\alpha,\beta)\). Since \(3\beta<\pi\), pairs at cyclic separation
\(k=1,2,3\) have shorter angle in \((k\alpha,k\beta)\). Their cosine
intervals have endpoints \(\cos(k\beta)\) and \(\cos(k\alpha)\),
which are rational by the multiple-angle formulas. These three
intervals and the near-antipodal fourth interval are disjoint.
Along with the inner and mixed pairs they yield six classes.

Let \(E_k^{\mathrm{out}}\) be the unordered outer edges of separation
\(k\). On each outer vertex the degrees of the four graphs are
\(2,2,2,1\). For any scalar radial function \(f\) and \(k=1,2,3\),
\begin{equation}
 \sum_{ij\in E_k^{\mathrm{out}}}\bigl(f(u_i)+f(u_j)\bigr)
 -2\sum_{ij\in E_4^{\mathrm{out}}}\bigl(f(u_i)+f(u_j)\bigr)=0.
                                                               \label{eq:cyclic}
\end{equation}
These identities permit class-specific radial terms in the upper
kernel while preserving its total over all pairs. At \(z^*\) the
contact cosines lie strictly inside their assigned intervals when the
Gram constraint is taken relative to its physical domain. The eight
first-class outer edges are adjacent in cyclic order. Indeed all outer
separations are at least \(\alpha\), and a nonadjacent edge must have
shorter angle at least \(2\alpha>\beta\). Thus the cyclic labels
remain consistent in a neighborhood of the candidate.

\section{A kernel with a summable bound}\label{sec:kernel}

For an angular order \(k\) use
\(H_k(u,v,t)=\real((z\bar w)^k)\). The recurrence
\(H_0=1\), \(H_1=t\), \(H_{k+1}=2tH_k-uvH_{k-1}\) makes each
\(H_k\) polynomial in the pair coordinates. Let \(r(u)\in\mathbb R^6\)
have the first three Chebyshev polynomials on the rescaled radial
interval containing \(u\) in its corresponding block, with zero
entries in the other block. The vector
\[
 m_k^*=\sum_i r(|z_i^*|^2)(z_i^*)^k
\]
is real by reflection symmetry of the candidate. For even \(k\) the
verifier chooses a nonzero coordinate \(p\) of \(m_k^*\) and uses
the five radial features
\(r_j(u)-r_p(u)(m_k^*)_j/(m_k^*)_p\), \(j\ne p\). Denote the resulting
vector by \(f_k(u)\). At odd \(k\) take \(f_k=r\), since antipodal
symmetry gives \(m_k^*=0\). In every case the candidate moment of
\(f_k(u)z^k\) is zero. For real symmetric positive definite matrices
\(A_k\), the kernel
\[
 K(z,w)=\sum_{k=0}^{11}H_k(u,v,t)\, f_k(u)^{\mathsf T}A_k f_k(v)
\]
sums over all ordered pairs to a sum of nonnegative Hermitian
quadratic forms in the corresponding moments. Add to \(K\) the
rank-one kernel \(2^{-40}q(u)^2q(v)^2\), where
\(q(u)=(u-a)(u-b)(u-c)\), and use the resulting \(K\) in
\begin{equation}
 h_0(z,w)=\frac{B}{45}-1+\frac59(u+v)-K(z,w)
                   -\frac{K(z,z)+K(w,w)}{18}.                 \label{eq:h0}
\end{equation}
Because \(\sum_{i<j}(u_i+u_j)=9Q=81\), we have
\(\sum_{i<j}h_{0,ij}=B-\tfrac12\sum_{i,j}K(z_i,z_j)\).

\subsection{Terms whose pair sums vanish}

Summing \eqref{eq:force}, and summing after multiplication by
\(z_i^2,z_i^3,z_i^4\), gives the identities
\begin{equation}
 \sum_i z_i=0,\qquad \sum_i u_i z_i=0,\qquad
 \sum_i (u_i-17/10)z_i^2=0,\qquad
 \sum_i (u_i-8/5)z_i^3=0.                                    \label{eq:ward}
\end{equation}
In these computations the sum of divided differences of \(z^3\)
over pairs is \((17/2)\sum_i z_i^2\), and for \(z^4\) it is
\(8\sum_i z_i^3\), using \(\sum_i z_i=0\). The mass identity
\(\sum_i(u_i-9/10)=0\) supplies a scalar condition.
For each identity use the zero moment and a corresponding
\(f_k\)-moment as factors of a real bilinear form, then distribute
the diagonal equally among the incident edges. The sum of the resulting
pair term vanishes. The \(f_k\)-moment at \(z^*\) also vanishes, so
the total of each such term has zero first derivative at the candidate.

The symmetrized row identities use the following expressions at
each fixed \(i\), summing over \(j\ne i\):
\[
 u_j-1+u_i/9,\qquad
 (8-\sigma_i)/9-\sigma_j,\qquad
 \frac{u_i-t_{ij}}{d_{ij}}-\frac{5u_i}{9}.
\]
Here \(\sigma_i\) is the indicator of the outer band. Their row sums
vanish by the mass, outer population, and real part of
\eqref{eq:force} multiplied by \(z_i\), respectively. The
multipliers of the mass rows use \(f_0\). Those of the force rows
use the band features multiplied by a polynomial that vanishes at
the candidate node or nodes in the band. The cyclic rows
\eqref{eq:cyclic} use the first three Chebyshev polynomials on \(O\).
The total of each row term has zero first derivative at the candidate.
For the force and mass rows this follows from the vanishing of both
factors after summing. Population and cyclic rows are identities in the
local band and cyclic labels.

\subsection{Second variations}

Take a real radial function \(g\) in the span of the six radial
features times their band-specific vanishing polynomials. Let
\(\xi_i=g(u_i)z_i\). For a pair, define
\[
\begin{aligned}
 N_g &=\tfrac59\bigl(ug(u)^2+vg(v)^2\bigr),\\
 T_g &=\frac{[g(u)(u-t)+g(v)(v-t)]^2
                 -(g(u)-g(v))^2(uv-t^2)}{d^2}.
\end{aligned}
\]
The real part of \(((\xi_i-\xi_j)/(z_i-z_j))^2\) equals \(T_g\).
Along the real lines with directions \(\xi\) and \(i\xi\),
the second derivatives of \(L\) are
\(-2\sum_{i<j}(N_g+T_g)\) and \(-2\sum_{i<j}(N_g-T_g)\).
Consequently the pair sums of the two quadratic forms are
nonnegative at a maximizer. Polarize both forms and contract
with real positive definite coefficient matrices. Subtracting
the resulting terms from \(h_0\) and including the zero-sum
corrections gives a class-dependent kernel \(h\) with
\begin{equation}
 \sum_{i<j}h_{ij}\ \le\ B-\tfrac12\sum_{i,j}K(z_i,z_j)
 \ \le\ B                                                     \label{eq:budget}
\end{equation}
at any global maximizer. The second-variation terms and their
first derivatives vanish at the candidate, because all six
displacement functions vanish there.

\section{Exact contact and continuous verification}\label{sec:verify}

\begin{lemma}\label{lem:verification}
With the coefficient data in \cite{data}, the kernel of
Section~\ref{sec:kernel} has positive definite matrices \(A_k\) and
positive definite second-variation coefficient matrices. In each
of the six necessary domains of Section~\ref{sec:domains},
\(h\ge\log d\). At every candidate pair, \(h=\log d\).
\end{lemma}

\begin{proof}[Certificate verification]
There are fifteen pair types in \eqref{eq:config}: one within the
inner pair, three in each rectangle, two between the inner pair
and each rectangle, and four between the two rectangles.
Equality of \(h\) and \(\log d\) at each type, with equality
of their first derivatives in the free directions, gives a
linear system with fifteen value rows and thirty-five derivative
rows. At an antipodal type with \(u=v\), exchange symmetry
identifies the two tangential derivatives on \(t=-\sqrt{uv}\).
At a same-radius interior type the radial derivatives are likewise
equal.

\emph{Exact coefficients.}
The system has six compatibility relationships. One is the
multiplicity-weighted sum of its value rows. The other five are
the derivatives of the weighted sum along the five parameters
of \eqref{eq:config}. For the candidate, the pair sum of \(h\)
is \(B\), and its first derivative is \(5\,dQ\), since the
moment penalties and all added terms have vanishing first
variation. Equation~\eqref{eq:force} gives \(d\log\Delta=5\,dQ\).
The five derivative relationships are independent, as the three
antipodal rows detect \(da,db,dc\), and within-rectangle changes
of the point products detect \(dx,dy\).
The seed consists of 312 numerical coefficients stored as rationals.
The exact correction is defined by solving a selected square
system of order 44 against the logarithms and derivatives at the
actual algebraic candidate. Directed dyadic interval elimination
at precision \(2^{-384}\) certifies its determinant is nonzero and
encloses the corrected coefficients. The six compatibility
relationships make the remaining rows exact. Interval \(LDL^{\mathsf T}\)
checks establish positivity of all twelve Gram matrices and both
second-variation matrices.

\emph{Contact neighborhoods.}
Write \(G=h-\log d\). At each of the twelve interior contact types,
strict positivity of the interval Hessian of \(G\), together with
exact vanishing of its constant and linear terms, proves \(G\ge0\)
in a neighborhood. For the three Gram-face contacts use the
positive definite Hessian of the restriction \(t=-\sqrt{uv}\)
and a positive inward derivative. The verifier establishes a
positive-radius neighborhood for each type.
Some types also use enlarged neighborhoods certified directly in
the cleared polynomial coordinates described below.

\emph{Interior and mixed domains.}
Put \(R=d^2h\) and \(F(d)=d^2\log d\). For each class, \(R\)
is polynomial in the pair variables. For \(d,A>0\),
\begin{equation}
 F(d)\le
  (\log A+3/2)d^2-2Ad+A^2/2+\frac{(d-A)^3}{3A},               \label{eq:taylor}
\end{equation}
since \(F^{(4)}(d)=-2/d^2<0\). In the mixed class and the first
three outer separation classes, the verifier starts with rational
boxes containing the domains, splitting at \(t=0\) where
needed. It rejects boxes disjoint from the necessary geometry.
For the remaining boxes it uses a certified contact neighborhood,
or takes \(A\) as the upper endpoint for \(d\), bounds
\(\log A\) above with an outward-rounded rational, and checks
that all tensor Bernstein coefficient lower bounds of \(R\)
minus the right side of \eqref{eq:taylor} are nonnegative.
An unaccepted box is divided into two. The validity of this
check follows from nonnegativity and the partition of unity
of the Bernstein basis. The verifier uses directed integer
arithmetic for conversions and subdivision. For enlarged
interior contact neighborhoods it takes \(A\) at the exact
contact distance and checks the Hessian of the polynomial
difference in coordinates centered at the contact.

\emph{The inner edge.}
The verifier first proves monotonicity toward increasing \(t\)
using \(d^3 G_t=dR_t+4R+2d^2\) on the domain swept by projection
to \(t=-\sqrt{uv}\). During this projection the initial
upper bound \(53/100\) for \(d\) can increase by at most
\(2(1-87/100)\sqrt{uv}\), and \(t\le-(87/100)\sqrt{uv}\)
throughout. With \(u=p^2,v=q^2,t=-pq\), the boundary has
\(d=(p+q)^2\). The verifier checks the resulting two-variable
polynomial against an upper Taylor bound for \(2s^4\log s\),
\(s=p+q\). A quadratic Taylor bound about the upper endpoint
of \(s\) applies since the third derivative
\(s(48\log s+52)\) is positive for \(s>1/2\). Substituting
an upper bound for the endpoint logarithm is valid since its
coefficient in the Taylor polynomial is positive.

\emph{The fourth outer separation.}
Write \(u=p^2,v=q^2,p\le q\) and use coordinates
\(s=p+q\), \(y'=q-p\), \(e=t+pq\).
Then \(0\le e\le pq/32\) and \(d=s^2-2e\).
With \(R\) expressed in these variables, \(R\) is even in \(y'\)
by exchange symmetry. On the axis \(y'=e=0\) the difference
\(R-2s^4\log s\) has a double zero at each of
\(s=2\sqrt b,2\sqrt c\). At each contact the verifier checks
positive axial curvature, nonnegativity of the polynomial
extension of \((R(s,y',0)-R(s,0,0))/(y')^2\), and
\(dR_e+4R+2d^2\ge0\) for the inward direction. These three
checks cover a box around each contact. The complement uses
Bernstein subdivision with a quadratic Taylor upper bound for
\(F\) about the upper endpoint of \(d\), since \(F'''(d)=2/d>0\).

\begin{center}
\begin{tabular}{lrr}
\toprule
Pair domain & Edge count & Subdivision nodes visited\\
\midrule
Inner edge & 1 & 1,122\\
Inner--outer & 16 & 5,030\\
Outer separation 1 & 8 & 11,743\\
Outer separation 2 & 8 & 890\\
Outer separation 3 & 8 & 30,653\\
Outer separation 4 & 4 & 829\\
\midrule
Total & 45 & 50,267\\
\bottomrule
\end{tabular}
\end{center}
The fixed report records the counts above and zero unresolved
boxes. Each accepted box has certified \(G\ge0\) throughout its
intersection with the necessary domain. Together with the
exact coefficient construction, this proves the lemma.
\end{proof}

\section{Sharpness and uniqueness}

At any global maximizer, Lemma~\ref{lem:verification} and
\eqref{eq:budget} imply
\[
 \log\Delta\le \sum_{i<j}h_{ij}\le B.
\]
The feasible candidate attains \(B\), so both inequalities
are equalities at every global maximizer. The positive
rank-one radial term forces \(\sum_i q(u_i)^2=0\).
All \(u_i\) therefore belong to \(\{a,b,c\}\).
Of the nonnegative integer triples summing to ten, only
\((2,4,4)\) has weighted mass nine, as follows by enumerating
the 66 triples using the isolating intervals for \(a,b,c\).

The first harmonic uses unprojected band features.
The inner constant and two independent outer linear features
extract a zero first moment on each radius. On the two-point
radius the points are antipodal. For four points on a circle
of squared radius \(r\), the elementary symmetric functions
satisfy \(e_3=e_4\overline{e_1}/r\). Thus \(e_1=0\) implies
that the root polynomial is even, and the four points are two
antipodal pairs.

For the second harmonic, vanishing of the projected moment
means the full moment is a complex multiple of \(m_2^*\).
The same radial interpolation now gives the second moment on
each of the three radii. The inner moment has modulus \(2a\),
so the complex multiple has modulus one. Rotate the whole
configuration to make it one. On each four-point radius, the
two squared images of the antipodal pairs have a fixed common
modulus and a prescribed nonzero sum, \(2x\) or \(2y\).
Two points on a circle with given nonzero sum have a fixed
midpoint and hence form a unique unordered pair. Taking the
two square roots of each of these images recovers
\eqref{eq:config}. Rescaling now proves Theorem~\ref{thm:main}.

\section*{Reproducibility}

The fixed certificate revision is
\href{\certbase/commit/\certrev}{\texttt{\certrev}}.
Its \certfile{docs/REPRODUCING.md}{replay guide} specifies the
dependency order and the source of the eight- and nine-point
results. From that checkout, the following commands replay the
ten-point dependency chain and the final check using Python's
standard library.
\begin{verbatim}
python -B -S -O scripts/verify_ten_point.py
python -B -S -O ten-point-sharp-contacts/verify_sharp.py
\end{verbatim}
The first command includes the second as its last stage. The
\certfile{ten-point-contact-geometry/root_box_jacobian.json}{Jacobian
result} supplies the fixed candidate-root data. The replay
independently checks the enclosed root and the exact arithmetic
certificates, without needing the numerical coefficient search
or a running Jacobian service. It checks its output against
the \certfile{ten-point-sharp-contacts/global_verification_report.json}
{fixed final report}.

\section*{Acknowledgments}

OpenAI Codex assisted with derivations, certificate implementation,
computation, and preparation of the manuscript. The candidate
root computation used the Jacobian toolkit \cite{jacobian}.

\begin{thebibliography}{9}
\bibitem{jxy}
Xiang Ji, Tiancheng Xu, and Wensheng Yu.
\newblock \emph{Optimal Configurations for Mordell's Problem}.
\newblock \href{https://arxiv.org/abs/2611.03088v1}{arXiv:2611.03088v1}, 2026.

\bibitem{n12}
Fengcan Jiang's Mordell research project.
\newblock \emph{The sharp twelve-point bound in Mordell's distance-product problem}.
\newblock \certfile{twelve-point-sharp-global/manuscript/README.md}{Manuscript and proof data}, 2026.

\bibitem{data}
Fengcan Jiang's Mordell research project.
\newblock \emph{Ten-point proof data, analytic reductions, and exact verifier}.
\newblock \certfile{ten-point-sharp-contacts/README.md}{Certificate archive}, 2026.

\bibitem{jacobian}
\emph{Jacobian}, composable tools for mathematical computation.
\newblock \href{https://github.com/morluto/jacobian}{Software repository}.
\end{thebibliography}
\end{document}
